\documentclass[11pt]{article}
\usepackage{metricsstyle}

\begin{document}

\lessonheader{2}{IV and GMM, the Gauss--Markov theorem, and the geometry of least squares}{}

\section{The linear model and its assumptions}

\[ y=X\beta+u,\qquad u\sim(0,\sigma^{2}I_n), \]
with $y$: $n\times1$, $X$: $n\times k$, $\beta$: $k\times1$ and $u$: $n\times1$. The
same model is carried under two alternative sets of assumptions about the
regressors and the errors, both keeping the error covariance spherical, plus a
rank condition:

\begin{flatlist}
  \item $X$ is fixed, \quad $u\sim N(0,\sigma^{2}I)$;
  \item $X$ is stochastic, \quad $u\sim(0,\sigma^{2}I)$;
  \item $\rank(X)=k$.
\end{flatlist}
The brace the notes draw around the error statements is labelled \emph{sphericity}:
one common variance $\sigma^{2}$, no correlation, which is what makes the algebra
below simple. The third line is the rank condition that makes $X'X$ invertible, so
that
\[ b=(X'X)^{-1}X'y \]
is defined.

\section{Orthogonality, and the need for instruments}

Together with $\rank(X)=k$, the two moment conditions that identify $\beta$ are
\[ \E(u)=0,\qquad \E(X'u)=0 \quad\text{(the \emph{orthogonality condition}).} \]
If the regressors are correlated with the error this fails. Then one looks for an
instrument matrix
\[ Z_{n\times p},\qquad p\ge k,\qquad \E(Z'u)=0, \]
that is, $p$ variables that are uncorrelated with $u$ but correlated with $X$
(formally, $\rank\E(Z'X)=k$). The moment condition $\E(Z'u)=0$ is what replaces $\E(X'u)=0$: it gives $p$ equations
for the $k$ unknown elements of $\beta$.

\section{The objective of an IV estimator}

The idea is to choose $\beta$ so that the sample moments $Z'u$ are as small as
possible. As written on the page,
\[ \min \|Z'u-\E(Z'u)\|,\qquad\text{which with } \E(Z'u)=0 \text{ is } \min\|Z'u\| . \]
Minimising a norm of a vector of moments is not the same as minimising the raw sum
of squares, because the moments do not carry the same variance. Weighting by the
inverse of their covariance (the metric of Lesson 1, $w'\Sigma^{-1}w$) gives
\[ \min (Z'u)'\bigl[\Var(Z'u)\bigr]^{-1}(Z'u). \]
Under sphericity $\Var(Z'u)=Z'\bigl(\sigma^{2}I_n\bigr)Z=\sigma^{2}Z'Z$, so the
inverse is $\sigma^{-2}(Z'Z)^{-1}$ and the objective becomes
\[ \frac{1}{\sigma^{2}}\,u'Z\,(Z'Z)^{-1}Z'u
   \;\Longrightarrow\;
   (y-X\beta)'Z\,(Z'Z)^{-1}Z'(y-X\beta), \]
the last step substituting $u=y-X\beta$. The factor $\sigma^{-2}$ is a constant and
drops out of the minimisation.

\section{The generalized IV (two-stage least squares) estimator}

Writing $P_z\equiv Z(Z'Z)^{-1}Z'$, minimise the objective over $b$:
\[
\begin{aligned}
  &\min_b\,(y-xb)'Z(Z'Z)^{-1}Z'(y-xb)\\
  &=\min_b\,(y-xb)'P_z(y-xb)\\
  &=\min_b\,\Bigl(y'P_zy-b'x'P_zy-y'P_zxb+b'x'P_zxb\Bigr).
\end{aligned}
\]
The two middle terms are the same scalar, so together they are $-2b'x'P_zy$, and the
first-order condition is
\[
\begin{aligned}
  -2x'P_zy+2x'P_zxb&=0\\
  x'P_zy&=(x'P_zx)\,b,
\end{aligned}
\]
hence
\[ b=(x'P_zx)^{-1}x'P_zy
   =(x'Z(Z'Z)^{-1}Z'x)^{-1}x'Z(Z'Z)^{-1}Z'y . \]
With more instruments than regressors, $p>k$ (the over-identified case), this is
the \emph{generalized IV} (two-stage least squares) estimator.

If instead there are exactly as many instruments as regressors, $p=k$, then $Z$ is
$n\times k$ and $Z'x$ is a square $k\times k$ matrix, invertible when the
instruments are relevant, $\rank(Z'x)=k$. Then
$(x'Z(Z'Z)^{-1}Z'x)^{-1}=(Z'x)^{-1}(Z'Z)(x'Z)^{-1}$, and the same formula collapses
to
\[ b=(Z'x)^{-1}(Z'Z)(x'Z)^{-1}\,x'Z(Z'Z)^{-1}Z'y=(Z'x)^{-1}Z'y, \]
the \emph{simple IV} estimator. With $p=k$ the choice of weights no longer matters:
$b$ sets the $k$ sample moments $Z'(y-xb)$ exactly to zero.

\section{Which estimator? A summary}

Reading the assumptions as a table:

\begin{flatlist}
  \item $y=X\beta+u$ with $\E(X'u)=0$ and $\Var(u)=\sigma^{2}I$: everything is
        spherical and the regressors are exogenous, so use least squares,
        $b=(x'x)^{-1}x'y$;
  \item $\E(X'u)\ne0$: the orthogonality condition fails, so bring in instruments
        $Z$ with $\E(Z'u)=0$ and use IV, $b_{iv}$; with $\Var(u)=\sigma^{2}I$ this is
        the estimator written out above;
  \item $\Var(u)\ne\sigma^{2}I_n$, i.e.\ $\Var(u)=\sigma^{2}\Omega_n$ with
        $\Omega\ne I_n$, while $X$ is fixed and $\E(x'u)=0$: the errors are
        non-spherical but exogenous, and the answer is \emph{generalized least
        squares};
  \item both failures at once, $\Var(u)\ne\sigma^{2}I_n$ \emph{and}
        $\E(x'u)\ne0$: generalized least squares is no longer consistent, and the
        answer is the \emph{generalized method of moments} (GMM).
\end{flatlist}

\section{The OLS estimator in the normal linear model}

Take the well-behaved case: $y=x\beta+u$ with $u\sim N(0,\sigma^{2}I)$, $X$ fixed and
$\rank(X)=k$, so that $x'x$ is invertible. The estimator is
\[
\begin{aligned}
  b&=(x'x)^{-1}x'y\\
   &=(x'x)^{-1}x'x\beta+(x'x)^{-1}x'u\\
   &=\beta+(x'x)^{-1}x'u .
\end{aligned}
\]
With $X$ fixed, $b$ is $\beta$ plus a linear function of $u$, so it inherits the
distribution of $u$: it is unbiased,
\[ \E(b)=\beta+(x'x)^{-1}x'\E(u)=\beta, \]
and its covariance is
\[
\begin{aligned}
  \Var(b)&=\E\bigl[(b-\beta)(b-\beta)'\bigr]
          =\E\bigl[(x'x)^{-1}x'uu'x(x'x)^{-1}\bigr]\\
         &=\sigma^{2}(x'x)^{-1},
\end{aligned}
\]
using $\E(uu')=\sigma^{2}I_n$, so the sampling distribution is normal,
\[ b\sim N\bigl(\beta,\ \sigma^{2}(x'x)^{-1}\bigr) . \]
The estimator is the \emph{BLUE} --- best linear unbiased estimator --- which is
the content of the Gauss--Markov theorem, proved next. The theorem uses only
$\E(u)=0$ and $\Var(u)=\sigma^{2}I$; normality is not needed for it, only for the
exact normal distribution of $b$.

\section{Gauss--Markov: no other linear unbiased estimator beats OLS}

Start from the model $y=X\beta+u$ and consider an arbitrary linear estimator
\[ b^{*}=Cy,\qquad C:\ k\times n,\quad y:\ n\times1, \]
just as least squares $b=(X'X)^{-1}X'y$ is $y$ premultiplied by the $k\times n$
matrix $(X'X)^{-1}X'$.
Every such $C$ can be written as the least-squares matrix plus a remainder,
\[ C=\bigl[(X'X)^{-1}X'+A\bigr], \]
so that
\[ b^{*}=Cy=\bigl[(X'X)^{-1}X'+A\bigr]y=(X'X)^{-1}X'y+Ay=b+Ay . \]
Substituting the model,
\[
\begin{aligned}
  b^{*}&=(x'x)^{-1}x'(x\beta+u)+A(x\beta+u)\\
       &=\beta+(x'x)^{-1}x'u+Ax\beta+Au .
\end{aligned}
\]
Taking expectations, $\E(u)=0$:
\[
\begin{aligned}
  \E(b^{*})&=\beta+(x'x)^{-1}x'\E(u)+Ax\beta+A\E(u)\\
           &=\beta+Ax\beta\;=\;\beta \qquad\text{when } Ax=0 .
\end{aligned}
\]
So $b^{*}$ is unbiased for every value of $\beta$ exactly when $A$ satisfies $Ax=0$. That same condition --- the
one that restricts attention to the class of unbiased linear estimators --- also
removes the cross terms below. For the variance,
\[
\begin{aligned}
  \Var(b^{*})&=\E\bigl[(b^{*}-\beta)(b^{*}-\beta)'\bigr]\\
             &=\E\Bigl[\bigl((x'x)^{-1}x'u+Au\bigr)\bigl(u'x(x'x)^{-1}+u'A'\bigr)\Bigr]\\
             &=\E\Bigl[(x'x)^{-1}x'uu'x(x'x)^{-1}+Auu'x(x'x)^{-1}
                    +(x'x)^{-1}x'uu'A'+Auu'A'\Bigr]\\
             &=\sigma^{2}(x'x)^{-1}+\sigma^{2}Ax(x'x)^{-1}
               +\sigma^{2}(x'x)^{-1}x'A'+\sigma^{2}AA' .
\end{aligned}
\]
With $Ax=0$ the two cross terms vanish and
\[
\begin{aligned}
  \Var(b^{*})&=\sigma^{2}(x'x)^{-1}+\sigma^{2}AA'\\
             &=\Var(b)+\sigma^{2}AA' .
\end{aligned}
\]
The extra term is $\sigma^{2}AA'$, which is positive semi-definite: for any
$c\in\R^{k}$, $c'AA'c=\|A'c\|^{2}\ge0$. (It is not in general positive definite;
it is nonzero whenever $A\ne0$, and zero only when $A=0$, i.e.\ $b^{*}=b$.) Adding
it can only make the covariance larger, so no linear unbiased estimator has a
smaller covariance than $b$: $\Var(b^{*})-\Var(b)$ is positive semi-definite, i.e.\
for every $c$, $\Var(c'b^{*})\ge\Var(c'b)$. Least squares is BLUE.

\section{Geometric interpretation}

In the model $y=X\beta+u$ the vector $y$ is a point of $\R^n$, and the columns of $X$ span a subspace
\[ \dim L(x)=k \]
only $k$-dimensional. Write the
least-squares decomposition as
\[ y=Xb+e, \]
where the fitted value $Xb=X(X'X)^{-1}X'y$ lies in that subspace and $e=y-Xb$ is
what is left over. The normal equations $X'(y-Xb)=0$ say exactly that $X'e=0$: the
residual is orthogonal to every column of $X$. (The true $X\beta$ also lies in the
subspace, but the error $u=y-X\beta$ is not orthogonal to it in general.)

\begin{figure}[ht]
\centering
\begin{tikzpicture}[scale=0.95]
  % the subspace spanned by the columns of X
  \draw[budget] (-1.4,0.2)
      .. controls (-1.1,1.4) and (0.2,2.0) .. (1.4,1.7)
      .. controls (2.6,1.4) and (3.5,0.9) .. (3.4,0.4)
      .. controls (3.3,-0.5) and (2.2,-1.4) .. (1.0,-1.5)
      .. controls (-0.2,-1.6) and (-1.2,-0.9) .. cycle;
  \node[ecnode] at (-0.35,1.05) {$L(x)$};
  % the origin
  \node[bundle] at (0,0) {};
  % y, and its projection
  \draw[ecaxis] (0,0) -- (4.2,2.2);
  \node[ecnode,anchor=west] at (4.28,2.24) {$y$};
  \draw[ecaxis] (0,0) -- (3.4,0.4);
  \node[ecnode,anchor=west] at (3.5,0.18) {$Xb$};
  % the residual, perpendicular to the subspace
  \draw[ecaxis] (3.4,0.4) -- (4.2,2.2);
  \node[ecnode,anchor=west] at (3.95,1.30) {$e$};
  % right angle at the foot of the perpendicular
  \draw[thin] (3.678,0.433) -- (3.792,0.689) -- (3.514,0.656);
  % annotation
  \node[ecnode] at (1.05,2.85) {orthogonal projection};
\end{tikzpicture}
\caption{$Xb$ is the orthogonal projection of $y$ onto the $k$-dimensional
subspace spanned by the columns of $X$, and the residual $e=y-Xb$ is orthogonal
to that subspace. Minimising $\|y-X\beta\|^{2}$ over $\beta$ is simply choosing the
point of the subspace closest to $y$.}
\end{figure}

\begin{transcriptnotes}
  \item The original opens with a title page reading only ``Lesson 2''.
  \item The three error statements are grouped by a hand-drawn brace labelled
        \emph{sphericity}; the arrangement is kept, the brace is not redrawn.
  \item The original mixes lower case ($x$) and upper case ($X$) for the regressor
        matrix, as in Lesson 1; the case is kept as written.
  \item $\E(Z'u)=0$ is used on the same page as the condition $\min\|Z'u-\E(Z'u)\|$.
        Both are transcribed as written, though with $\E(Z'u)=0$ the second term of
        the norm vanishes; the norm is then replaced by
        $(Z'u)'[\Var(Z'u)]^{-1}Z'u$, which is what makes the weights appear.
  \item The constant $\sigma^{-2}$ is carried into the last objective line here; the
        original drops it there, since it does not affect the minimiser.
  \item \textbf{Corrected:} the first-order condition was written
        $-2b'x'P_zy+2x'P_zxb=0$; the derivative of $-2b'x'P_zy$ with respect to $b$
        is $-2x'P_zy$ (no $b'$), which is what the next line
        $x'P_zy=(x'P_zx)b$ uses.
  \item \textbf{Corrected:} in the just-identified collapse the middle line reads
        $(Z'x)^{-1}(Z'Z)\,x'(x'Z)^{-1}x'Z(Z'Z)^{-1}Z'y$ with a stray $x'$ after
        $(Z'Z)$, and the line above has $x'Z'(Z'Z)^{-1}$ for $x'Z(Z'Z)^{-1}$; both
        are typeset without the slip.
  \item The page annotates the generalized-IV formula with ``$p>k$'' and the
        simple-IV case with ``$K=p$, $Z$ is $n\times K$''; the annotation is kept
        as ``over-identified''/``just-identified''.
  \item $b\sim(\beta,\sigma^{2}(x'x)^{-1})$ is written on the page as the shorthand
        $(\text{mean},\text{covariance})$; since $u$ is normal the distribution is
        normal, and it is typeset $N(\beta,\sigma^{2}(x'x)^{-1})$.
  \item \textbf{Corrected:} the expectation of $b^{*}$ was written
        $\E(b)+(x'x)^{-1}x'\E(u)+\dots$; $\E(b)$ already contains the term
        $(x'x)^{-1}x'\E(u)$, so the first term is $\beta$, not $\E(b)$.
  \item \textbf{Corrected:} ``$\sigma^{2}AA'$ is positive definite when $A\ne0$''
        is not true in general ($A$ is $k\times n$ and $AA'$ is positive definite
        only if $A$ has full row rank $k$); the conclusion only needs positive
        semi-definiteness, which always holds.
  \item \textbf{Added:} the geometry page contains only $y=X\beta+u$,
        $y\in\R^n$ (written ``$y=\R^n$''), $\dim\ell(x)=k$, the sketch and the
        words ``orthogonal projection''. The decomposition $y=Xb+e$, the normal
        equations $X'e=0$ and the remark on $u$ are added explanation, as are the
        labels $Xb$, $e$ and $y$ in the figure (the sketch labels only the
        subspace).
  \item The just-identified case is written with $K$ on the page; it is typeset
        $p=k$ to match the dimensions used everywhere else ($k$ and $K$ are the
        same number of regressors). The page states
        $\rank(Z)=K$; what the formula needs is $\rank(Z'x)=k$ (relevance), and
        the step showing the collapse to $(Z'x)^{-1}Z'y$ is added.
  \item On the last page the subspace is labelled $\ell(x)$ in the sketch (a cursive
        lower-case letter); it is the column space of $X$ and is typeset $L(x)$.
  \item The sketch also has an arrow from the origin to the label $\ell(x)$
        inside the subspace; it is read as a pointer to the label and not redrawn.
  \item The Gauss--Markov page writes ``$\Var(b^{*})=\sigma^{2}(x'x)^{-1}+\sigma^{2}AA'$
        where $A=$'' and leaves the condition unfinished; it is $Ax=0$, stated
        two lines earlier on the page.
  \item \textbf{Corrected:} the conclusion was written ``for any,
        $\Var(b^{*})>\Var(b)$''; the correct statement is that
        $\Var(b^{*})-\Var(b)$ is positive semi-definite ($\ge$, with equality
        when $A=0$).
  \item \textbf{Added:} the remarks that the moments are weighted by the inverse
        of their covariance (with the link to Lesson~1), that relevance means
        $\rank\E(Z'X)=k$, that $\sigma^{-2}$ drops out of the minimisation, that
        with $p=k$ the sample moments are set exactly to zero, that GLS is no longer
        consistent when $\E(x'u)\ne0$, and that Gauss--Markov does not need
        normality, are explanation, not on the page.
  \item Pages of the original, for tracing back: page 2 the model and the
        assumptions, the IV moment condition and the IV objective; page 3 the two-stage-least-squares and simple-IV formulas;
        page 4 the summary of assumptions and estimators and the start of
        $b=\beta+(x'x)^{-1}x'u$; page 5 $\E(b)$,
        $\Var(b)$ and the start of $b^{*}$; page 6 the Gauss--Markov proof;
        page 7 the geometric interpretation.
\end{transcriptnotes}

\end{document}
